\begingroup
\ArithmeticHead{278}{L}{ANALYSIS. NACHLASS.}
\SourceLeaf{285b}
\pdfbookmark[1]{Erläuterungen zu [VI.] und [VII.]}{x1-agm-erlaeuterung-vi-vii}
\begin{center}\textls[120]{Erläuterungen zu [VI.] und [VII.], S. 207—212..}\end{center}
\par\indent Die Aufzeichnung der Scheda Ae im Abschnitt [VI.] kann als Vorarbeit zu der Pars altera der Abhandlung von 1800 (Werke III, S. 372) angesehen werden und bildet zugleich durch die Ausführung einiger in jener Abhandlung nur angedeuteten Rechnungen eine Ergänzung zu ihr. In den Gleichungen (1), (2) des art. 2. ist
\[x'=\tfrac12(x+y),\quad y'=\sqrt{xy},\]
\noindent ebenso bedeuten weiterhin $x'',y''$ u.s.w. die folgenden Elemente im Algorithmus des agM. In der Formel S. 208, Z. 5 fehlen in der Handschrift gewisse Faktoren; die Formel muß lauten
\[\begin{aligned}
\frac{dM(x,y)}{dy}=\frac12\frac M{x'}\Bigl(&1+2\frac{y'(x'x'-y'y')}{y(xx-yy)}\,.\,\frac{x'}{x''}\\
&+4\frac{y''(x''x''-y''y'')}{y(xx-yy)}\,.\,\frac{x'}{x'''}+\ldots\Bigr).
\end{aligned}\]
\par\indent Die Aufzeichnung der Scheda Af, Abschnitt [VII.], geht sehr erheblich über den Standpunkt der Abhandlung von 1800 (Werke III, S. 361) hinaus, indem sie bereits alle wesentlichen Bestandteile für eine selbständige Theorie der \textls[120]{Modulfunktion} enthält. Sie gewinnt besondere historische Bedeutung dadurch, daß sie Ergebnisse, die sonst nur aus erheblich später zu datierenden Aufzeichnungen bekannt waren, für das Jahr 1801 sicherstellt. So die Formel [1], die übersichtliche Zusammenstellung der Ausdrücke von $a,b,\sqrt{a^2-b^2}$ durch die Reihen $p,q,r$ am Anfang des art. [3.], und die wichtige Relation [10], die hier anscheinend unabhängig von der Differentialgleichung [12] aus der Darstellung [2] hergeleitet worden ist.
\clearpage
\fancyhead{}\fancyhead[R]{\small279}\fancyhead[C]{\scriptsize BEMERKUNGEN ZUM AGM. ABSCHNITT VIII.}
\SourceLeaf{286a}
\noindent Man vergleiche art. [1.] des Abschnitts [VIII.], S. 213, die artt. [2.] und [4.] des Abschnitts [IX.], S. 218 und 220, ferner Werke III, S. 468, 469, die alle aus sehr viel späterer Zeit stammen. In Gleichung [3], S. 210, ist
\[0{,}1660958116=\log_{10}\frac2{m\pi},\]
\noindent wo $m$ der Modul der \textsc{Brigg}schen Logarithmen ist; in [5] beziehungsweise [7] ist
\[f=\frac1{M(a,b)},\quad F=\frac1{M(A,B)}.\]
\par\indent Besonders hervorzuheben ist noch das \textit{Theorema} im art. [5.]; dieses spielt nämlich eine wichtige Rolle bei der im Handbuch 16, Bb, S. 111—112 gegebenen Lösung des Umkehrproblems des elliptischen Integrals erster Gattung (abgedruckt im Auszug Werke III, S. 401, 402, vollständig hier als Abschnitt [II.] des folgenden Hauptstücks). \textsc{Gauss} führt es daselbst (siehe Werke III, S. 401) an mit der Bemerkung »wovon jedoch der Beweis tiefer liegt«.
\clearpage
\endgroup

